% Options for packages loaded elsewhere
% Options for packages loaded elsewhere
\PassOptionsToPackage{unicode}{hyperref}
\PassOptionsToPackage{hyphens}{url}
\PassOptionsToPackage{dvipsnames,svgnames,x11names}{xcolor}
%
\documentclass[
  letterpaper,
  oneside,
  open=any]{scrbook}
\usepackage{xcolor}
\usepackage[top=25mm,left=25mm,right=25mm,bottom=30mm]{geometry}
\usepackage{amsmath,amssymb}
\setcounter{secnumdepth}{5}
\usepackage{iftex}
\ifPDFTeX
  \usepackage[T1]{fontenc}
  \usepackage[utf8]{inputenc}
  \usepackage{textcomp} % provide euro and other symbols
\else % if luatex or xetex
  \usepackage{unicode-math} % this also loads fontspec
  \defaultfontfeatures{Scale=MatchLowercase}
  \defaultfontfeatures[\rmfamily]{Ligatures=TeX,Scale=1}
\fi
\usepackage{lmodern}
\ifPDFTeX\else
  % xetex/luatex font selection
\fi
% Use upquote if available, for straight quotes in verbatim environments
\IfFileExists{upquote.sty}{\usepackage{upquote}}{}
\IfFileExists{microtype.sty}{% use microtype if available
  \usepackage[]{microtype}
  \UseMicrotypeSet[protrusion]{basicmath} % disable protrusion for tt fonts
}{}
\makeatletter
\@ifundefined{KOMAClassName}{% if non-KOMA class
  \IfFileExists{parskip.sty}{%
    \usepackage{parskip}
  }{% else
    \setlength{\parindent}{0pt}
    \setlength{\parskip}{6pt plus 2pt minus 1pt}}
}{% if KOMA class
  \KOMAoptions{parskip=half}}
\makeatother
% Make \paragraph and \subparagraph free-standing
\makeatletter
\ifx\paragraph\undefined\else
  \let\oldparagraph\paragraph
  \renewcommand{\paragraph}{
    \@ifstar
      \xxxParagraphStar
      \xxxParagraphNoStar
  }
  \newcommand{\xxxParagraphStar}[1]{\oldparagraph*{#1}\mbox{}}
  \newcommand{\xxxParagraphNoStar}[1]{\oldparagraph{#1}\mbox{}}
\fi
\ifx\subparagraph\undefined\else
  \let\oldsubparagraph\subparagraph
  \renewcommand{\subparagraph}{
    \@ifstar
      \xxxSubParagraphStar
      \xxxSubParagraphNoStar
  }
  \newcommand{\xxxSubParagraphStar}[1]{\oldsubparagraph*{#1}\mbox{}}
  \newcommand{\xxxSubParagraphNoStar}[1]{\oldsubparagraph{#1}\mbox{}}
\fi
\makeatother


\usepackage{longtable,booktabs,array}
\newcounter{none} % for unnumbered tables
\usepackage{calc} % for calculating minipage widths
% Correct order of tables after \paragraph or \subparagraph
\usepackage{etoolbox}
\makeatletter
\patchcmd\longtable{\par}{\if@noskipsec\mbox{}\fi\par}{}{}
\makeatother
% Allow footnotes in longtable head/foot
\IfFileExists{footnotehyper.sty}{\usepackage{footnotehyper}}{\usepackage{footnote}}
\makesavenoteenv{longtable}
\usepackage{graphicx}
\makeatletter
\newsavebox\pandoc@box
\newcommand*\pandocbounded[1]{% scales image to fit in text height/width
  \sbox\pandoc@box{#1}%
  \Gscale@div\@tempa{\textheight}{\dimexpr\ht\pandoc@box+\dp\pandoc@box\relax}%
  \Gscale@div\@tempb{\linewidth}{\wd\pandoc@box}%
  \ifdim\@tempb\p@<\@tempa\p@\let\@tempa\@tempb\fi% select the smaller of both
  \ifdim\@tempa\p@<\p@\scalebox{\@tempa}{\usebox\pandoc@box}%
  \else\usebox{\pandoc@box}%
  \fi%
}
% Set default figure placement to htbp
\def\fps@figure{htbp}
\makeatother





\setlength{\emergencystretch}{3em} % prevent overfull lines

\providecommand{\tightlist}{%
  \setlength{\itemsep}{0pt}\setlength{\parskip}{0pt}}



 


\makeatletter
\@ifpackageloaded{tcolorbox}{}{\usepackage[skins,breakable]{tcolorbox}}
\@ifpackageloaded{fontawesome5}{}{\usepackage{fontawesome5}}
\definecolor{quarto-callout-color}{HTML}{909090}
\definecolor{quarto-callout-note-color}{HTML}{0758E5}
\definecolor{quarto-callout-important-color}{HTML}{CC1914}
\definecolor{quarto-callout-warning-color}{HTML}{EB9113}
\definecolor{quarto-callout-tip-color}{HTML}{00A047}
\definecolor{quarto-callout-caution-color}{HTML}{FC5300}
\definecolor{quarto-callout-color-frame}{HTML}{acacac}
\definecolor{quarto-callout-note-color-frame}{HTML}{4582ec}
\definecolor{quarto-callout-important-color-frame}{HTML}{d9534f}
\definecolor{quarto-callout-warning-color-frame}{HTML}{f0ad4e}
\definecolor{quarto-callout-tip-color-frame}{HTML}{02b875}
\definecolor{quarto-callout-caution-color-frame}{HTML}{fd7e14}
\makeatother
\makeatletter
\@ifpackageloaded{bookmark}{}{\usepackage{bookmark}}
\makeatother
\makeatletter
\@ifpackageloaded{caption}{}{\usepackage{caption}}
\AtBeginDocument{%
\ifdefined\contentsname
  \renewcommand*\contentsname{Table of contents}
\else
  \newcommand\contentsname{Table of contents}
\fi
\ifdefined\listfigurename
  \renewcommand*\listfigurename{List of Figures}
\else
  \newcommand\listfigurename{List of Figures}
\fi
\ifdefined\listtablename
  \renewcommand*\listtablename{List of Tables}
\else
  \newcommand\listtablename{List of Tables}
\fi
\ifdefined\figurename
  \renewcommand*\figurename{Figure}
\else
  \newcommand\figurename{Figure}
\fi
\ifdefined\tablename
  \renewcommand*\tablename{Table}
\else
  \newcommand\tablename{Table}
\fi
}
\@ifpackageloaded{float}{}{\usepackage{float}}
\floatstyle{ruled}
\@ifundefined{c@chapter}{\newfloat{codelisting}{h}{lop}}{\newfloat{codelisting}{h}{lop}[chapter]}
\floatname{codelisting}{Listing}
\newcommand*\listoflistings{\listof{codelisting}{List of Listings}}
\makeatother
\makeatletter
\makeatother
\makeatletter
\@ifpackageloaded{caption}{}{\usepackage{caption}}
\@ifpackageloaded{subcaption}{}{\usepackage{subcaption}}
\makeatother
\usepackage{bookmark}
\IfFileExists{xurl.sty}{\usepackage{xurl}}{} % add URL line breaks if available
\urlstyle{same}
\hypersetup{
  pdftitle={BMS CHO Cell Productivity Knowledge Graph},
  pdfauthor={Elucidata × BMS},
  colorlinks=true,
  linkcolor={\#6d28d9},
  filecolor={Maroon},
  citecolor={Blue},
  urlcolor={\#ea580c},
  pdfcreator={LaTeX via pandoc}}


\title{BMS CHO Cell Productivity Knowledge Graph}
\usepackage{etoolbox}
\makeatletter
\providecommand{\subtitle}[1]{% add subtitle to \maketitle
  \apptocmd{\@title}{\par {\large #1 \par}}{}{}
}
\makeatother
\subtitle{User Manual --- Version 3}
\author{Elucidata × BMS}
\date{March 16, 2026}
\begin{document}
\frontmatter
\maketitle

\renewcommand*\contentsname{Table of contents}
{
\setcounter{tocdepth}{1}
\tableofcontents
}

\mainmatter
\bookmarksetup{startatroot}

\chapter{Welcome}\label{welcome}

\section{About this Manual}\label{about-this-manual}

This manual is your end-to-end guide for working with the \textbf{BMS
CHO Cell Productivity Knowledge Graph (KG)}---a multi-omics,
multi-evidence platform built to accelerate the identification of
genetic engineering targets that improve CHO cell productivity.

The KG integrates:

\begin{itemize}
\tightlist
\item
  \textbf{BMS in-house data} --- transcriptomics, proteomics, and
  metabolomics from Study 1 and Study 2 alongside bioprocess readouts
  (qP, titre, IVCD, lactate, cell density)
\item
  \textbf{CHO-specific public omics data} --- 13 datasets from 10+
  publications comparing high vs.~low producers
\item
  \textbf{Causal literature evidence} --- cause-effect relationships
  curated from \textasciitilde500 genetic perturbation publications
\item
  \textbf{External databases} --- NCBI, UniProt, GO, KEGG, ChEMBL,
  PubMed, Cellosaurus
\end{itemize}

The bioinformatics analyses embedded in the KG include differential
expression / abundance (DEG/DAA), WGCNA co-expression networks, MOFA
multi-omics integration, pathway enrichment (GSEA \& ORA), and
transcription-factor enrichment.

\begin{center}\rule{0.5\linewidth}{0.5pt}\end{center}

\section{How to Use this Manual}\label{how-to-use-this-manual}

{\def\LTcaptype{none} % do not increment counter
\begin{longtable}[]{@{}
  >{\raggedright\arraybackslash}p{(\linewidth - 2\tabcolsep) * \real{0.5000}}
  >{\raggedright\arraybackslash}p{(\linewidth - 2\tabcolsep) * \real{0.5000}}@{}}
\toprule\noalign{}
\begin{minipage}[b]{\linewidth}\raggedright
If you want to\ldots{}
\end{minipage} & \begin{minipage}[b]{\linewidth}\raggedright
Go to\ldots{}
\end{minipage} \\
\midrule\noalign{}
\endhead
\bottomrule\noalign{}
\endlastfoot
Understand what the KG contains before querying &
Section~\ref{sec-orientation} \\
Browse internal BMS omics data and run DEG / WGCNA / MOFA queries &
Section~\ref{sec-internal} \\
Look up what published studies say about a gene &
Section~\ref{sec-literature} \\
Check whether BMS data agrees with published findings &
Section~\ref{sec-cross-evidence} \\
Explore pathway enrichment results & Section~\ref{sec-pathways} \\
Identify master regulators / transcription factors &
Section~\ref{sec-tf} \\
Test a specific biological hypothesis & Section~\ref{sec-hypothesis} \\
Shortlist high-confidence engineering targets &
Section~\ref{sec-targets} \\
Search, export, or understand data coverage &
Section~\ref{sec-utilities} \\
\end{longtable}
}

\begin{center}\rule{0.5\linewidth}{0.5pt}\end{center}

\section{Manual Structure}\label{manual-structure}

This manual is organised into \textbf{nine chapters} covering all major
use-cases, followed by a glossary of key terms.

\begin{center}\rule{0.5\linewidth}{0.5pt}\end{center}

\section{Key Conventions}\label{key-conventions}

\begin{tcolorbox}[enhanced jigsaw, arc=.35mm, bottomrule=.15mm, bottomtitle=1mm, breakable, colback=white, colbacktitle=quarto-callout-note-color!10!white, colframe=quarto-callout-note-color-frame, coltitle=black, left=2mm, leftrule=.75mm, opacityback=0, opacitybacktitle=0.6, rightrule=.15mm, title=\textcolor{quarto-callout-note-color}{\faInfo}\hspace{0.5em}{Note}, titlerule=0mm, toprule=.15mm, toptitle=1mm]

\textbf{Query boxes} throughout this manual show verbatim
natural-language queries (NLQ) you can type directly into the KG
interface. Results may vary slightly depending on your cohort or filter
settings.

\end{tcolorbox}

\begin{tcolorbox}[enhanced jigsaw, arc=.35mm, bottomrule=.15mm, bottomtitle=1mm, breakable, colback=white, colbacktitle=quarto-callout-tip-color!10!white, colframe=quarto-callout-tip-color-frame, coltitle=black, left=2mm, leftrule=.75mm, opacityback=0, opacitybacktitle=0.6, rightrule=.15mm, title=\textcolor{quarto-callout-tip-color}{\faLightbulb}\hspace{0.5em}{Tip}, titlerule=0mm, toprule=.15mm, toptitle=1mm]

\textbf{Evidence badges} tag information by analytical layer:

{DEG/DAA} {WGCNA} {MOFA} {Literature} {GEM (in progress)}

\end{tcolorbox}

\begin{center}\rule{0.5\linewidth}{0.5pt}\end{center}

\emph{For support, contact
\href{mailto:yogesh.lakhotia@elucidata.io}{\nolinkurl{yogesh.lakhotia@elucidata.io}}}

\bookmarksetup{startatroot}

\chapter{Orientation}\label{orientation}

What the KG is, what it contains, and how to talk to it

\hfill\break

\section{Overview}\label{sec-orientation}

This chapter is your starting point. Before running any query, spend a
few minutes here to understand how the KG is structured---you will get
significantly better results.

\begin{center}\rule{0.5\linewidth}{0.5pt}\end{center}

\section{What is the Knowledge Graph?}\label{sec-kg-intro}

A \textbf{Knowledge Graph (KG)} is a structured, semantic representation
of data that connects real-world entities---genes, proteins,
metabolites, pathways, samples, publications---through typed
relationships.

The BMS CHO Productivity KG uses a \textbf{property graph} model (Neo4j)
with three primitives:

{\def\LTcaptype{none} % do not increment counter
\begin{longtable}[]{@{}
  >{\raggedright\arraybackslash}p{(\linewidth - 4\tabcolsep) * \real{0.3333}}
  >{\raggedright\arraybackslash}p{(\linewidth - 4\tabcolsep) * \real{0.3333}}
  >{\raggedright\arraybackslash}p{(\linewidth - 4\tabcolsep) * \real{0.3333}}@{}}
\toprule\noalign{}
\begin{minipage}[b]{\linewidth}\raggedright
Primitive
\end{minipage} & \begin{minipage}[b]{\linewidth}\raggedright
Description
\end{minipage} & \begin{minipage}[b]{\linewidth}\raggedright
Examples
\end{minipage} \\
\midrule\noalign{}
\endhead
\bottomrule\noalign{}
\endlastfoot
\textbf{Nodes} (Entities) & Objects or concepts & Gene, Protein,
Metabolite, Sample, Pathway, Factor, Cohort, Study \\
\textbf{Edges} (Connections) & Typed relationships &
\texttt{gene\_expressed\_in}, \texttt{upregulated},
\texttt{gene\_part\_of}, \texttt{tf\_regulates\_cohort\_degs} \\
\textbf{Properties} (Attributes) & Metadata on nodes/edges & log2FC,
adj. p-value, kME, MOFA weight, Spearman r \\
\end{longtable}
}

The KG is queried through a natural-language interface that converts
your question into Cypher graph queries automatically. You can also use
the NLQ tips in Section~\ref{sec-query-tips}.

\begin{center}\rule{0.5\linewidth}{0.5pt}\end{center}

\section{Data Model Overview}\label{data-model-overview}

The complete data model connects the following entity types:

Key relationships to know:

\begin{itemize}
\tightlist
\item
  \textbf{Gene → Sample}: \texttt{gene\_expressed\_in} (TPM values, per
  sample)
\item
  \textbf{Gene → Cohort}: \texttt{upregulated} / \texttt{downregulated}
  / \texttt{non\_regulated} (DESeq2 results)
\item
  \textbf{Gene / Protein / Metabolite → Network}:
  \texttt{gene\_part\_of} (WGCNA module membership, kME)
\item
  \textbf{Gene / Protein / Metabolite → Factor}:
  \texttt{gene\_contributes\_to\_factor} (MOFA weights)
\item
  \textbf{Gene → Positive/Negative Phenotypes}: Spearman correlation
  with qP, titre, IVCD
\item
  \textbf{Gene → Cohort}: \texttt{tf\_regulates\_cohort\_degs} (TF
  enrichment)
\item
  \textbf{Pathway → Cohort / Network / Factor}:
  \texttt{kegg\_pathway\_enriched\_in} / \texttt{go\_term\_enriched\_in}
\end{itemize}

\begin{center}\rule{0.5\linewidth}{0.5pt}\end{center}

\section{What Data Lives in the KG?}\label{sec-data-contents}

\subsection{In-House BMS Data}\label{in-house-bms-data}

\textbf{Study 1} --- multi-omics time-series (Days 3, 5, 9, 11, 14):

{\def\LTcaptype{none} % do not increment counter
\begin{longtable}[]{@{}lllll@{}}
\toprule\noalign{}
Molecule & Clones & High qP & Low qP & Omics \\
\midrule\noalign{}
\endhead
\bottomrule\noalign{}
\endlastfoot
C144 & 5 & 2 & 3 & Transcriptomics, Proteomics, Metabolomics \\
Nivo & 3 & 1 & 2 & Transcriptomics, Proteomics, Metabolomics \\
\end{longtable}
}

\textbf{Study 2} --- proteomics only (Days 5, 7, 9, 11):

{\def\LTcaptype{none} % do not increment counter
\begin{longtable}[]{@{}llll@{}}
\toprule\noalign{}
Molecule & Clones & High qP & Low qP \\
\midrule\noalign{}
\endhead
\bottomrule\noalign{}
\endlastfoot
Milvaxin & 6 & 3 & 3 \\
COVID2 & 6 & 3 & 3 \\
Nivo & 3 & 1 & 2 \\
FXIaR & 6 & 3 & 3 \\
\end{longtable}
}

Bioprocess measurements in Study 1: \textbf{qP, titre, IVCD, lactate,
O₂, cell density}, and others.

\subsection{Public Data}\label{public-data}

\begin{itemize}
\tightlist
\item
  \textbf{11 transcriptomics datasets} from 10 publications (Study 1 +
  public), comparing high vs.~low or high vs.~non-producing clones
\item
  \textbf{3 proteomics datasets} (Studies 1 \& 2 + public)
\item
  \textbf{1 metabolomics dataset} (Study 1)
\item
  \textbf{\textasciitilde500 publications} on genetic perturbation →
  productivity outcomes (45 ingested at the time of writing)
\end{itemize}

\subsection{External Databases}\label{external-databases}

NCBI Gene, UniProt, ChEBI, KEGG (pathways, reactions, enzymes,
orthologs), Gene Ontology (BP, MF, CC), PubMed, Cellosaurus.

\begin{center}\rule{0.5\linewidth}{0.5pt}\end{center}

\section{Cohort Naming Conventions}\label{sec-cohorts}

The KG stores results for \textbf{883 cohorts} across four categories.
When querying for DEG or enrichment results, \textbf{always name your
cohort explicitly} for best retrieval.

{\def\LTcaptype{none} % do not increment counter
\begin{longtable}[]{@{}
  >{\raggedright\arraybackslash}p{(\linewidth - 4\tabcolsep) * \real{0.3333}}
  >{\raggedright\arraybackslash}p{(\linewidth - 4\tabcolsep) * \real{0.3333}}
  >{\raggedright\arraybackslash}p{(\linewidth - 4\tabcolsep) * \real{0.3333}}@{}}
\toprule\noalign{}
\begin{minipage}[b]{\linewidth}\raggedright
Category
\end{minipage} & \begin{minipage}[b]{\linewidth}\raggedright
Count
\end{minipage} & \begin{minipage}[b]{\linewidth}\raggedright
Example
\end{minipage} \\
\midrule\noalign{}
\endhead
\bottomrule\noalign{}
\endlastfoot
Productivity-Based (High vs.~Low, all clones \& days) & 6 &
\texttt{Nivo\_study1\_all\_high\_vs\_Nivo\_study1\_all\_low} \\
Clone-Wise (across all days) & 36 &
\texttt{Nivo\_clone8\_study1\_all\_vs\_Nivo\_clone49\_study1\_all} \\
Temporal (same clone, different days) & 190 &
\texttt{Nivo\_clone49\_day11\_study1\_vs\_Nivo\_clone49\_day8\_study1} \\
Temporal Cross-Clone & 630 &
\texttt{Nivo\_clone8\_day11\_study1\_vs\_Nivo\_clone49\_day8\_study1} \\
Public (high vs.~low / non-producer) & 19 &
\texttt{GSE295486\_high\_producer\_vs\_GSE295486\_non\_producer} \\
\end{longtable}
}

\begin{tcolorbox}[enhanced jigsaw, arc=.35mm, bottomrule=.15mm, bottomtitle=1mm, breakable, colback=white, colbacktitle=quarto-callout-tip-color!10!white, colframe=quarto-callout-tip-color-frame, coltitle=black, left=2mm, leftrule=.75mm, opacityback=0, opacitybacktitle=0.6, rightrule=.15mm, title=\textcolor{quarto-callout-tip-color}{\faLightbulb}\hspace{0.5em}{Tip}, titlerule=0mm, toprule=.15mm, toptitle=1mm]

Prefer \textbf{Productivity-Based} cohorts when asking global ``which
genes drive high productivity?'' questions. Use Clone-Wise or Temporal
cohorts for clone-specific or time-resolved investigations.

\end{tcolorbox}

\begin{center}\rule{0.5\linewidth}{0.5pt}\end{center}

\section{How to Phrase Natural-Language Queries}\label{sec-query-tips}

The KG interface accepts plain English. Follow these principles for
reliable results:

\begin{enumerate}
\def\labelenumi{\arabic{enumi}.}
\tightlist
\item
  \textbf{Name the entity type} --- say ``gene'', ``protein'', or
  ``metabolite'', not just a symbol.
\item
  \textbf{Name the cohort or condition} --- include molecule (Nivo /
  C144) and study when relevant.
\item
  \textbf{State the direction} --- ``upregulated in high producers'',
  ``negatively correlated with qP''.
\item
  \textbf{State the evidence layer} --- ``in WGCNA module M1'', ``with
  MOFA Factor 1 weight \textgreater{} 0.3''.
\item
  \textbf{Add thresholds explicitly} --- ``adj. p-value \textless{} 0.05
  and \textbar log2FC\textbar{} \textgreater{} 2''.
\end{enumerate}

Show me genes upregulated in high vs.~low producers in Nivo Study 1 with
adj. p-value \textless{} 0.05 and absolute log2 fold change greater than
2.

Which proteins are positively correlated with qP (Spearman r
\textgreater{} 0.5) in Study 1?

What KEGG pathways are enriched in WGCNA Module M1 for the Nivo high
vs.~low productivity cohort?

\begin{center}\rule{0.5\linewidth}{0.5pt}\end{center}

\section{Understanding Query Results}\label{understanding-query-results}

Results are returned as \textbf{tables}, \textbf{network views}, or
\textbf{charts} depending on the query type:

\begin{itemize}
\tightlist
\item
  \textbf{Tables} --- for DEG lists, correlation rankings, module
  membership
\item
  \textbf{Network views} --- for TF-gene regulatory graphs, WGCNA
  modules, KG subgraphs
\item
  \textbf{Bar / violin charts} --- for pathway enrichment scores,
  correlation distributions
\end{itemize}

Key columns you will encounter:

{\def\LTcaptype{none} % do not increment counter
\begin{longtable}[]{@{}
  >{\raggedright\arraybackslash}p{(\linewidth - 2\tabcolsep) * \real{0.5000}}
  >{\raggedright\arraybackslash}p{(\linewidth - 2\tabcolsep) * \real{0.5000}}@{}}
\toprule\noalign{}
\begin{minipage}[b]{\linewidth}\raggedright
Column
\end{minipage} & \begin{minipage}[b]{\linewidth}\raggedright
Meaning
\end{minipage} \\
\midrule\noalign{}
\endhead
\bottomrule\noalign{}
\endlastfoot
\texttt{log2FC} & Log₂ fold change (DEG); positive = higher in first
group \\
\texttt{adj\_pvalue} & Benjamini-Hochberg corrected p-value \\
\texttt{kME} & Module eigengene connectivity (WGCNA hub score; closer to
1 = stronger hub) \\
\texttt{gs\_qp} / \texttt{gs\_titre} / \texttt{gs\_ivcd} & Gene
significance scores (WGCNA correlation with trait) \\
\texttt{mofa\_weight} & MOFA feature weight; sign indicates direction,
magnitude indicates importance \\
\texttt{spearman\_r} & Spearman correlation of feature abundance with
productivity phenotype \\
\texttt{NES} & Normalised enrichment score (GSEA); positive = enriched
in first group \\
\end{longtable}
}

\bookmarksetup{startatroot}

\chapter{Internal Data Exploration}\label{internal-data-exploration}

Browsing BMS omics data, differential expression, co-expression
networks, and multi-omics integration

\hfill\break

\section{Overview}\label{sec-internal}

This chapter covers everything you can learn from \textbf{BMS in-house
data alone}. It is organised into four progressive layers of analysis,
each building on the previous:

\begin{center}\rule{0.5\linewidth}{0.5pt}\end{center}

\section{A · Data Browsing}\label{sec-browsing}

\subsection{Samples, Clones, and
Timepoints}\label{samples-clones-and-timepoints}

What samples are available for the Nivo molecule in Study 1?

Which days have both transcriptomics and proteomics measurements in
Study 1?

The KG stores sample-level coverage. Study 1 has multi-omics samples at
\textbf{Days 3, 5, 9, 11, and 14}. Study 2 proteomics spans \textbf{Days
5, 7, 9, and 11}.

\subsection{Productivity Measurements Across
Clones}\label{productivity-measurements-across-clones}

Show qP, titre, and IVCD values across all Nivo clones in Study 1.

\subsection{Molecule-Level Filtering}\label{molecule-level-filtering}

Use the molecule filter to restrict results to a single drug candidate:

\begin{itemize}
\tightlist
\item
  \texttt{Nivo} --- nivolumab (anti-PD-1 mAb)
\item
  \texttt{C144} --- C144 antibody
\end{itemize}

Show me differentially expressed genes for C144 high vs.~low producers
only.

\begin{center}\rule{0.5\linewidth}{0.5pt}\end{center}

\section{B · Differential Expression \& Abundance}\label{sec-deg}

\subsection{Method Reference}\label{method-reference}

{\def\LTcaptype{none} % do not increment counter
\begin{longtable}[]{@{}llll@{}}
\toprule\noalign{}
Omics Layer & Method & Covariate & Normalisation \\
\midrule\noalign{}
\endhead
\bottomrule\noalign{}
\endlastfoot
Transcriptomics (DEG) & DESeq2 & Time & TPM \\
Proteomics (DAA) & T-test & Time & Median central log₂ \\
Metabolomics (DAA) & T-test & Time & Median central log₂ \\
\end{longtable}
}

Results are stored per \textbf{cohort} as directed edges from features
to cohort nodes with properties \texttt{log2FC}, \texttt{adj\_pvalue},
and \texttt{regulation} (upregulated / downregulated / non\_regulated).

\subsection{Finding Upregulated Genes in High
Producers}\label{finding-upregulated-genes-in-high-producers}

Which genes are upregulated in high vs.~low Nivo producers in Study 1
with adj. p-value \textless{} 0.05 and \textbar log2FC\textbar{}
\textgreater{} 2?

\subsection{Cross-Cohort Consistency}\label{cross-cohort-consistency}

A gene is considered a \textbf{robust signal} if it is consistently
differentially expressed across ≥ 7/10 high-vs-low cohorts. The
shortlisting funnel applied to Study 1 data:

{\def\LTcaptype{none} % do not increment counter
\begin{longtable}[]{@{}ll@{}}
\toprule\noalign{}
Filter step & Genes remaining \\
\midrule\noalign{}
\endhead
\bottomrule\noalign{}
\endlastfoot
All DEGs across 232 cohorts & 11,481 \\
adj. p-value \textless{} 0.05 \& & log2FC \\
Directional consistency (RNA \& protein both up or both down) & 775 \\
Upregulated in high producers & 330 \\
Concordant RNA--protein \& Spearman r \textgreater{} 0.5 with qP &
\textbf{63} \\
\end{longtable}
}

Show genes consistently upregulated in at least 7 out of 10 high vs.~low
producer cohorts in Study 1.

\subsection{Comparing Across Cohorts}\label{comparing-across-cohorts}

Compare DEG results between the C144 and Nivo high-vs-low cohorts. Which
genes are shared?

\subsection{Protein and Metabolite Differential
Abundance}\label{protein-and-metabolite-differential-abundance}

Which proteins are differentially abundant in high vs.~low Nivo
producers in Study 1?

Show differentially abundant metabolites in the pellet fraction for high
vs.~low C144 producers.

\begin{center}\rule{0.5\linewidth}{0.5pt}\end{center}

\section{C · Co-expression Networks (WGCNA)}\label{sec-wgcna}

\subsection{What is WGCNA?}\label{what-is-wgcna}

\textbf{Weighted Gene Co-expression Network Analysis (WGCNA)} clusters
features that vary together across samples into \textbf{modules}. Each
module is summarised by its \textbf{Module Eigengene (ME)}---the first
principal component of all features in that module---which is then
correlated with bioprocess traits (qP, titre, IVCD).

\begin{quote}
\textbf{Published precedent:} Melville et al.~(2011, \emph{Journal of
Biotechnology}) applied WGCNA to 295 microarrays from 121 CHO cultures
and identified six transcriptional modules, five of which correlated
with bioprocess variables including qP. The BMS KG replicates and
extends this approach across all three omics layers.
\end{quote}

\subsection{Finding Productivity-Correlated
Modules}\label{finding-productivity-correlated-modules}

Which WGCNA gene modules are positively correlated with qP in Study 1?

Show all WGCNA networks (gene, protein, and metabolite) with significant
correlation to titre.

Key modules identified:

{\def\LTcaptype{none} % do not increment counter
\begin{longtable}[]{@{}llll@{}}
\toprule\noalign{}
Module & Omics & ME--qP correlation & Biological theme \\
\midrule\noalign{}
\endhead
\bottomrule\noalign{}
\endlastfoot
M1 (Pro) & Gene & Positive & ER protein processing / UPR \\
M18 (Anti) & Gene & Negative & Cell cycle / DNA replication \\
\end{longtable}
}

\subsection{Exploring Module
Membership}\label{exploring-module-membership}

List all genes in WGCNA module M1 with their kME scores.

\subsection{Hub Genes}\label{hub-genes}

Hub genes have the highest \textbf{kME} (module eigengene connectivity)
and are most representative of the module's biological program.

What are the top 10 hub genes in WGCNA module M1?

\subsection{Gene Significance Scores}\label{gene-significance-scores}

Gene significance (GS) scores quantify the correlation between each
gene's expression and a phenotype. Available as \texttt{gs\_qp},
\texttt{gs\_ivcd}, and \texttt{gs\_titre}.

Show genes in module M1 with gs\_qp \textgreater{} 0.5, ranked by
descending kME.

\subsection{WGCNA for Proteins and
Metabolites}\label{wgcna-for-proteins-and-metabolites}

The same WGCNA workflow was applied to proteomics and metabolomics data,
yielding \textbf{5 protein modules} and \textbf{5 metabolite modules}.

Which protein WGCNA modules correlate with IVCD in Study 1?

Show metabolite module members with significant correlation to qP,
ranked by gs\_qp.

\begin{center}\rule{0.5\linewidth}{0.5pt}\end{center}

\section{D · Multi-Omics Factor Analysis (MOFA)}\label{sec-mofa}

\subsection{What is MOFA?}\label{what-is-mofa}

\textbf{Multi-Omics Factor Analysis (MOFA)} is an unsupervised
dimensionality-reduction method that integrates transcriptomics,
proteomics, and metabolomics simultaneously to uncover \textbf{latent
factors}---hidden biological programs that explain variation across
samples.

\subsection{Significant Factors and Their Biological
Meaning}\label{significant-factors-and-their-biological-meaning}

{\def\LTcaptype{none} % do not increment counter
\begin{longtable}[]{@{}
  >{\raggedright\arraybackslash}p{(\linewidth - 4\tabcolsep) * \real{0.3333}}
  >{\raggedright\arraybackslash}p{(\linewidth - 4\tabcolsep) * \real{0.3333}}
  >{\raggedright\arraybackslash}p{(\linewidth - 4\tabcolsep) * \real{0.3333}}@{}}
\toprule\noalign{}
\begin{minipage}[b]{\linewidth}\raggedright
Factor
\end{minipage} & \begin{minipage}[b]{\linewidth}\raggedright
Correlation pattern
\end{minipage} & \begin{minipage}[b]{\linewidth}\raggedright
Interpretation
\end{minipage} \\
\midrule\noalign{}
\endhead
\bottomrule\noalign{}
\endlastfoot
\textbf{Factor 1} & qP +0.848 · IVCD +0.860 · Titre +0.904 &
Full-spectrum productivity signal --- higher cell efficiency, volumetric
yield, and cumulative biomass \\
\textbf{Factor 2} & qP positive · IVCD negative &
High-efficiency-per-cell phenotype --- not driven by biomass
accumulation \\
\textbf{Factor 3} & Proteomics-only & Post-translational regulatory axis
decoupled from transcription \\
\textbf{Factor 4} & Transcriptomics-only & Transcriptional program not
yet reflected in protein or metabolites at sampled time points \\
\textbf{Factor 5} & Titre \& qP strongly negative & Pure metabolic
suppressor acting on secretory machinery \\
\end{longtable}
}

\begin{tcolorbox}[enhanced jigsaw, arc=.35mm, bottomrule=.15mm, bottomtitle=1mm, breakable, colback=white, colbacktitle=quarto-callout-important-color!10!white, colframe=quarto-callout-important-color-frame, coltitle=black, left=2mm, leftrule=.75mm, opacityback=0, opacitybacktitle=0.6, rightrule=.15mm, title=\textcolor{quarto-callout-important-color}{\faExclamation}\hspace{0.5em}{Important}, titlerule=0mm, toprule=.15mm, toptitle=1mm]

\textbf{Factor 1} is the most actionable factor. Genes / proteins /
metabolites with high positive Factor 1 weights are pro-productivity
candidates; those with high negative weights are anti-productivity
candidates.

\end{tcolorbox}

\subsection{Querying MOFA Results}\label{querying-mofa-results}

Which genes have the highest positive MOFA Factor 1 weights?

Show proteins with MOFA Factor 1 weight above 0.3.

Which metabolites contribute to Factor 5? What pathways are enriched?

Show per-sample MOFA Factor 1 scores and overlay with qP measurements.

\subsection{Cross-Omics Consistency: DEG × WGCNA ×
MOFA}\label{cross-omics-consistency-deg-wgcna-mofa}

The overlap of all three evidence layers produces the highest-confidence
candidates:

\begin{itemize}
\tightlist
\item
  \textbf{Pro-productivity triple overlap}: 20 genes (upregulated in
  ≥7/10 cohorts AND in WGCNA M1 AND positive Factor 1 weight)
\item
  \textbf{Anti-productivity triple overlap}: 34 genes
\end{itemize}

Show me genes that are consistently upregulated in high producers (≥7/10
cohorts), members of a pro-productivity WGCNA module, and have positive
MOFA Factor 1 weight.

Enriched pathways in the triple-overlap candidates:

{\def\LTcaptype{none} % do not increment counter
\begin{longtable}[]{@{}lll@{}}
\toprule\noalign{}
Direction & Pathway & Significance \\
\midrule\noalign{}
\endhead
\bottomrule\noalign{}
\endlastfoot
Pro-productivity & Protein processing in ER & ★★★ \\
Pro-productivity & Cell adhesion molecules & ★★★ \\
Pro-productivity & Lysosome & ★★★ \\
Anti-productivity & Cell cycle & ★★★ \\
Anti-productivity & DNA replication & ★★★ \\
Anti-productivity & Homologous recombination & ★★★ \\
\end{longtable}
}

\bookmarksetup{startatroot}

\chapter{Public Literature Evidence}\label{public-literature-evidence}

What published studies say about genes, proteins, and productivity

\hfill\break

\section{Overview}\label{sec-literature}

The KG integrates \textbf{cause-effect relationships} curated from
\textasciitilde500 publications on genetic perturbation in the context
of CHO productivity. This chapter covers how to query that evidence
layer.

\begin{center}\rule{0.5\linewidth}{0.5pt}\end{center}

\section{How Literature Evidence is
Structured}\label{how-literature-evidence-is-structured}

Publications were scouted from \textbf{\textasciitilde45,000 CHO-related
papers} through a two-stage filter:

\begin{enumerate}
\def\labelenumi{\arabic{enumi}.}
\tightlist
\item
  \textbf{Filter 1} --- grouped into three categories:

  \begin{itemize}
  \tightlist
  \item
    \textbf{(A)} Gene/protein/metabolite target manipulation and cell
    engineering studies ✓ \emph{included}
  \item
    \textbf{(B)} Bioprocess optimisation
  \item
    \textbf{(C)} Other studies
  \end{itemize}
\item
  \textbf{Filter 2} --- Inclusion: productivity insights through target
  manipulation or omics-driven cell-line engineering. Exclusion: not
  directly relevant to target manipulation impacting productivity.
\end{enumerate}

Result: \textbf{\textasciitilde500 publications} describing genetic
interventions (overexpression, knockdown, knockout) and their effects on
productivity.

Each extracted relationship is stored as a \textbf{directed edge} with
properties:

{\def\LTcaptype{none} % do not increment counter
\begin{longtable}[]{@{}ll@{}}
\toprule\noalign{}
Property & Description \\
\midrule\noalign{}
\endhead
\bottomrule\noalign{}
\endlastfoot
\texttt{cause} & Gene / protein / metabolite manipulated \\
\texttt{manipulation} & OE (overexpression), KD (knockdown), KO
(knockout) \\
\texttt{effect\_entity} & Productivity phenotype affected \\
\texttt{effect\_direction} & Positive (+) or negative (−) \\
\texttt{cell\_line} & e.g., CHO-K1 GS KO \\
\texttt{target\_product} & e.g., IgG1 \\
\texttt{reactor} & e.g., Fed-batch \\
\texttt{pmid} & PubMed ID of source publication \\
\texttt{evidence\_sentence} & Source text from the paper \\
\end{longtable}
}

\begin{center}\rule{0.5\linewidth}{0.5pt}\end{center}

\section{Key Effect Entities}\label{key-effect-entities}

The KG tracks effects across four biological dimensions:

\textbf{Productivity:} specific\_productivity\_qp · product\_titer ·
volumetric\_productivity · secretion\_rate · product\_yield ·
production\_stability

\textbf{Growth:} viable\_cell\_density\_vcd · cell\_viability · IVCD ·
specific\_growth\_rate · apoptosis\_rate ·
cell\_cycle\_phase\_distribution · cloning\_efficiency

\textbf{Metabolic:} lactate\_production/consumption ·
ammonia\_production · glucose\_consumption · glutamine\_metabolism ·
redox\_balance · metabolic\_flux · lipid\_metabolism

\textbf{Quality:} glycosylation\_profile · fucosylation · sialylation ·
product\_purity · protein\_fragmentation · aggregation\_level ·
host\_cell\_protein\_content

\begin{center}\rule{0.5\linewidth}{0.5pt}\end{center}

\section{Querying Literature
Evidence}\label{querying-literature-evidence}

\subsection{Finding All Genes Linked to
Productivity}\label{finding-all-genes-linked-to-productivity}

Which genes are linked to specific productivity (qP) in published CHO
studies?

Show all genes with published evidence for improving product titer in
CHO cells.

\subsection{Looking Up a Specific
Gene}\label{looking-up-a-specific-gene}

What does the literature say about HSPA5 (GRP78/BiP) and CHO
productivity?

Summarise all published evidence for ST14 in CHO cells, including
manipulation type and effect.

\subsection{Filtering by Phenotype}\label{filtering-by-phenotype}

Which genes, when overexpressed, positively affect product titer in CHO?

Show genes whose knockout reduces protein fragmentation.

\subsection{Filtering by Source
Organism}\label{filtering-by-source-organism}

Show only CHO-specific literature evidence for HSPA5.

Are there human or mouse studies for FOXO1 that might be relevant to CHO
productivity?

\subsection{Retrieving Supporting Evidence
Sentences}\label{retrieving-supporting-evidence-sentences}

Show the supporting sentences from publications linking HSPA5
overexpression to increased titre.

\subsection{Genes with Experimental
Manipulation}\label{genes-with-experimental-manipulation}

Which genes have been experimentally overexpressed in CHO cells and
shown a positive effect on qP? Show PMID and evidence.

\begin{center}\rule{0.5\linewidth}{0.5pt}\end{center}

\section{Worked Example: HSPA5
(GRP78/BiP)}\label{worked-example-hspa5-grp78bip}

{Literature}

HSPA5 encodes the ER chaperone GRP78/BiP. Published manipulation studies
support overexpression as a productivity lever:

{\def\LTcaptype{none} % do not increment counter
\begin{longtable}[]{@{}llll@{}}
\toprule\noalign{}
PMID & Manipulation & Effect & Product \\
\midrule\noalign{}
\endhead
\bottomrule\noalign{}
\endlastfoot
32624757 & GRP78 OE & ↑ Titre & mAb \\
24081924 & BiP co-expression & ↑ qP & Recombinant protein \\
1373378 & GRP78 OE & ↑ Productivity & IgG \\
\end{longtable}
}

\textbf{Mechanism:} UPR activation expands ER folding capacity →
improves secretory efficiency.

\begin{center}\rule{0.5\linewidth}{0.5pt}\end{center}

\section{Worked Example: ST14}\label{worked-example-st14}

{Literature}

ST14 encodes Matriptase, a serine protease. Knockout studies show
improved product quality:

{\def\LTcaptype{none} % do not increment counter
\begin{longtable}[]{@{}lll@{}}
\toprule\noalign{}
PMID & Manipulation & Effect \\
\midrule\noalign{}
\endhead
\bottomrule\noalign{}
\endlastfoot
29777593 & ST14 KO & ↓ Protein fragmentation \\
29777593 & ST14 KO & ↑ Product purity \\
\end{longtable}
}

\textbf{Mechanism:} Matriptase degrades secreted product → knockout
eliminates fragmentation and improves purity.

\bookmarksetup{startatroot}

\chapter{Cross-Evidence Integration}\label{cross-evidence-integration}

Concordance and discordance between BMS internal data and published
literature

\hfill\break

\section{Overview}\label{sec-cross-evidence}

The highest-confidence engineering candidates are those supported by
\textbf{both} BMS internal data and independent published literature.
This chapter explains how to find, interpret, and resolve agreement---or
disagreement---between these two evidence streams.

\begin{center}\rule{0.5\linewidth}{0.5pt}\end{center}

\section{Why Cross-Evidence Analysis
Matters}\label{why-cross-evidence-analysis-matters}

Internal data alone can surface many differentially expressed
genes---but not all are validated targets. Literature alone reflects
findings in other cell lines, products, and contexts. When both sources
\textbf{agree in direction}, confidence in the target rises
substantially.

\begin{tcolorbox}[enhanced jigsaw, arc=.35mm, bottomrule=.15mm, bottomtitle=1mm, breakable, colback=white, colbacktitle=quarto-callout-note-color!10!white, colframe=quarto-callout-note-color-frame, coltitle=black, left=2mm, leftrule=.75mm, opacityback=0, opacitybacktitle=0.6, rightrule=.15mm, title=\textcolor{quarto-callout-note-color}{\faInfo}\hspace{0.5em}{Note}, titlerule=0mm, toprule=.15mm, toptitle=1mm]

The KG stores BMS omics results and literature cause-effect
relationships as separate edge types on the same gene nodes. This allows
a single query to traverse both layers simultaneously.

\end{tcolorbox}

\begin{center}\rule{0.5\linewidth}{0.5pt}\end{center}

\section{Checking if a Literature Gene Appears in BMS
Data}\label{checking-if-a-literature-gene-appears-in-bms-data}

Is HSPA5 differentially expressed in BMS Study 1 high vs.~low producer
cohorts?

Is HSPA5 a member of a WGCNA module correlated with productivity in
Study 1?

Does HSPA5 have a significant MOFA Factor 1 weight in Study 1?

\begin{center}\rule{0.5\linewidth}{0.5pt}\end{center}

\section{Direction Concordance}\label{direction-concordance}

The most important check: does the \textbf{direction} of the BMS signal
agree with the published manipulation effect?

{\def\LTcaptype{none} % do not increment counter
\begin{longtable}[]{@{}
  >{\raggedright\arraybackslash}p{(\linewidth - 4\tabcolsep) * \real{0.3333}}
  >{\raggedright\arraybackslash}p{(\linewidth - 4\tabcolsep) * \real{0.3333}}
  >{\raggedright\arraybackslash}p{(\linewidth - 4\tabcolsep) * \real{0.3333}}@{}}
\toprule\noalign{}
\begin{minipage}[b]{\linewidth}\raggedright
BMS signal
\end{minipage} & \begin{minipage}[b]{\linewidth}\raggedright
Published effect
\end{minipage} & \begin{minipage}[b]{\linewidth}\raggedright
Verdict
\end{minipage} \\
\midrule\noalign{}
\endhead
\bottomrule\noalign{}
\endlastfoot
Gene upregulated in high producers & Overexpression → ↑ qP & ✅
Concordant --- strong OE candidate \\
Gene downregulated in high producers & Knockout → ↑ titer & ✅
Concordant --- strong KO candidate \\
Gene upregulated in high producers & Overexpression → ↓ qP & ⚠️
Discordant --- investigate context \\
Gene not DE in BMS & Overexpression → ↑ qP & 🔍 BMS data silent ---
check other layers \\
\end{longtable}
}

For genes with published overexpression evidence showing positive qP
effect, which are also upregulated in high producers in BMS Study 1?

\begin{center}\rule{0.5\linewidth}{0.5pt}\end{center}

\section{Explaining Discordance}\label{explaining-discordance}

When BMS and literature disagree, common explanations include:

\begin{itemize}
\tightlist
\item
  \textbf{Source organism}: human or mouse study; CHO physiology may
  differ
\item
  \textbf{Product type}: a gene relevant for mAb production may behave
  differently for fusion proteins
\item
  \textbf{Cell line background}: CHO-K1 vs.~CHO-DG44 vs.~CHO-S show
  regulatory differences
\item
  \textbf{Timepoint}: literature perturbation may have been measured at
  a different culture stage
\end{itemize}

The literature shows positive productivity effect for gene X, but BMS
data show it is downregulated in high producers. Are the literature
studies from CHO cells or human cells?

\begin{center}\rule{0.5\linewidth}{0.5pt}\end{center}

\section{Finding Genes with Multi-Omics + Literature
Support}\label{finding-genes-with-multi-omics-literature-support}

The most actionable candidates:

Show genes that are (1) consistently upregulated in high producers in
BMS Study 1 across ≥7 cohorts, (2) members of a pro-productivity WGCNA
module, (3) have positive MOFA Factor 1 weight, AND (4) have published
evidence linking them to increased qP or titer.

The intersection of these criteria converges to a small, high-confidence
set. For example, the HSPA5 / ST14 pair shown in
Section~\ref{sec-literature} each pass all four filters.

\begin{center}\rule{0.5\linewidth}{0.5pt}\end{center}

\section{Genes Unique to BMS Data}\label{genes-unique-to-bms-data}

Some BMS signals have \textbf{no published precedent}---these are
potentially novel targets.

Show genes consistently upregulated in high producers in BMS Study 1
that are NOT mentioned in any ingested publication.

These novel signals may be CHO-cell-line-specific or molecule-specific
effects not yet studied in the field.

\begin{center}\rule{0.5\linewidth}{0.5pt}\end{center}

\section{Mapping BMS Data to Public
Omics}\label{mapping-bms-data-to-public-omics}

The KG includes 13 public omics datasets from external CHO studies. You
can check whether a BMS-identified gene is also differentially expressed
in independent public datasets.

Is HSPA5 differentially expressed in any of the public high-vs-low
producer omics datasets?

Which genes from the BMS Study 1 triple-overlap shortlist are also DE in
at least 3 public CHO datasets?

\bookmarksetup{startatroot}

\chapter{Pathway \& Functional
Enrichment}\label{pathway-functional-enrichment}

KEGG, GO, GSEA, and ORA results stored in the KG

\hfill\break

\section{Overview}\label{sec-pathways}

Pathway enrichment results are pre-computed and \textbf{stored directly
in the KG} as edges between pathway nodes and cohort, network, or factor
nodes. You do not need to run enrichment yourself---you query the
results.

\begin{center}\rule{0.5\linewidth}{0.5pt}\end{center}

\section{Enrichment Methods}\label{enrichment-methods}

{\def\LTcaptype{none} % do not increment counter
\begin{longtable}[]{@{}
  >{\raggedright\arraybackslash}p{(\linewidth - 6\tabcolsep) * \real{0.2500}}
  >{\raggedright\arraybackslash}p{(\linewidth - 6\tabcolsep) * \real{0.2500}}
  >{\raggedright\arraybackslash}p{(\linewidth - 6\tabcolsep) * \real{0.2500}}
  >{\raggedright\arraybackslash}p{(\linewidth - 6\tabcolsep) * \real{0.2500}}@{}}
\toprule\noalign{}
\begin{minipage}[b]{\linewidth}\raggedright
Method
\end{minipage} & \begin{minipage}[b]{\linewidth}\raggedright
Full name
\end{minipage} & \begin{minipage}[b]{\linewidth}\raggedright
Input
\end{minipage} & \begin{minipage}[b]{\linewidth}\raggedright
What it finds
\end{minipage} \\
\midrule\noalign{}
\endhead
\bottomrule\noalign{}
\endlastfoot
\textbf{GSEA} & Gene Set Enrichment Analysis & Ranked gene list (by
log2FC or MOFA weight) & Pathways with coordinated shifts across the
full ranked list \\
\textbf{ORA} & Over-Representation Analysis & Binary gene list (DEGs or
module members) & Pathways with more hits than expected by chance \\
\end{longtable}
}

Both methods are applied to:

\begin{itemize}
\tightlist
\item
  \textbf{DEG cohorts} (per molecule, per condition)
\item
  \textbf{WGCNA modules} (significant modules only)
\item
  \textbf{MOFA factors} (top 80th percentile features per factor per
  omics layer)
\end{itemize}

Databases used: \textbf{KEGG Pathways}, \textbf{Gene Ontology Biological
Process (GO:BP)}, \textbf{Gene Ontology Molecular Function (GO:MF)},
\textbf{Gene Ontology Cellular Component (GO:CC)}.

\begin{center}\rule{0.5\linewidth}{0.5pt}\end{center}

\section{Querying Pathway Enrichment}\label{querying-pathway-enrichment}

\subsection{On DEG Cohorts}\label{on-deg-cohorts}

Which KEGG pathways are enriched in high vs.~low Nivo producers in Study
1?

Show GO Biological Process terms enriched in the downregulated genes of
C144 high vs.~low producers.

Filter enrichment results for the Nivo high-vs-low cohort to GSEA
results only.

\subsection{On WGCNA Modules}\label{on-wgcna-modules}

What KEGG pathways are enriched in WGCNA gene module M1?

Show GO terms enriched in the anti-productivity WGCNA gene module M18.

\subsection{On MOFA Factors}\label{on-mofa-factors}

Which pathways are enriched in the top-weighted genes for MOFA Factor 1?

Show pathway enrichment results for MOFA Factor 5, separated by positive
and negative weight genes.

\begin{center}\rule{0.5\linewidth}{0.5pt}\end{center}

\section{Key Biological Themes}\label{key-biological-themes}

Results from DEG, WGCNA, and MOFA all converge on consistent pathway
themes:

\subsection{Pro-Productivity Pathways}\label{pro-productivity-pathways}

{\def\LTcaptype{none} % do not increment counter
\begin{longtable}[]{@{}
  >{\raggedright\arraybackslash}p{(\linewidth - 4\tabcolsep) * \real{0.3333}}
  >{\raggedright\arraybackslash}p{(\linewidth - 4\tabcolsep) * \real{0.3333}}
  >{\raggedright\arraybackslash}p{(\linewidth - 4\tabcolsep) * \real{0.3333}}@{}}
\toprule\noalign{}
\begin{minipage}[b]{\linewidth}\raggedright
Pathway
\end{minipage} & \begin{minipage}[b]{\linewidth}\raggedright
Source
\end{minipage} & \begin{minipage}[b]{\linewidth}\raggedright
NES / Significance
\end{minipage} \\
\midrule\noalign{}
\endhead
\bottomrule\noalign{}
\endlastfoot
Protein processing in endoplasmic reticulum & MOFA F1 GSEA & +3.3 ★★★ \\
Cell adhesion molecules & MOFA F1 GSEA & +3.1 ★★★ \\
Lysosome & MOFA F1 GSEA & +3.0 ★★★ \\
Protein export & WGCNA M1 ORA & Significant \\
Antigen processing and presentation & WGCNA M1 ORA & Significant \\
\end{longtable}
}

\subsection{Anti-Productivity
Pathways}\label{anti-productivity-pathways}

{\def\LTcaptype{none} % do not increment counter
\begin{longtable}[]{@{}lll@{}}
\toprule\noalign{}
Pathway & Source & NES / Significance \\
\midrule\noalign{}
\endhead
\bottomrule\noalign{}
\endlastfoot
Cell cycle & MOFA F1 GSEA & −3.2 ★★★ \\
DNA replication & MOFA F1 GSEA & −3.0 ★★★ \\
Homologous recombination & MOFA F1 GSEA & −2.9 ★★★ \\
Fanconi anemia pathway & MOFA F1 GSEA & −2.7 ★★★ \\
Mismatch repair & MOFA F1 GSEA & −2.6 ★★★ \\
\end{longtable}
}

\begin{tcolorbox}[enhanced jigsaw, arc=.35mm, bottomrule=.15mm, bottomtitle=1mm, breakable, colback=white, colbacktitle=quarto-callout-tip-color!10!white, colframe=quarto-callout-tip-color-frame, coltitle=black, left=2mm, leftrule=.75mm, opacityback=0, opacitybacktitle=0.6, rightrule=.15mm, title=\textcolor{quarto-callout-tip-color}{\faLightbulb}\hspace{0.5em}{Tip}, titlerule=0mm, toprule=.15mm, toptitle=1mm]

\textbf{Biological interpretation:} High producers allocate resources to
ER protein folding and secretion (UPR / secretory pathway upregulated)
while suppressing cell proliferation programs (cell cycle / DNA
replication downregulated). This is consistent with the canonical
productivity vs.~growth trade-off in CHO cells.

\end{tcolorbox}

\begin{center}\rule{0.5\linewidth}{0.5pt}\end{center}

\section{Comparing Enriched Pathways Across
Molecules}\label{comparing-enriched-pathways-across-molecules}

Compare enriched KEGG pathways between Nivo and C144 high-vs-low
producer cohorts. Which pathways are shared, and which are
molecule-specific?

\begin{center}\rule{0.5\linewidth}{0.5pt}\end{center}

\section{KEGG Metabolic Network
Traversal}\label{kegg-metabolic-network-traversal}

Beyond pathway enrichment, the KG stores the full KEGG metabolic network
(enzymes, reactions, substrates, products), enabling mechanistic
queries.

Which KEGG pathway does gene LDHA belong to?

What enzymatic activity does LDHA have, and which metabolic reactions
does it catalyse?

Find all genes in the glycolysis / gluconeogenesis KEGG pathway that are
differentially expressed in high vs.~low Nivo producers.

Trace LDHA to its substrate and product: what metabolites are connected?

Identify metabolic bottlenecks by combining DEG and KEGG data for the
lactate metabolism pathway.

\bookmarksetup{startatroot}

\chapter{Transcription Factor
Analysis}\label{transcription-factor-analysis}

Master regulators of high-producer gene expression programs

\hfill\break

\section{Overview}\label{sec-tf}

Transcription factor (TF) enrichment analysis identifies upstream
regulators whose target genes are over-represented in a DEG list or
WGCNA module. This chapter covers how to find, filter, and interpret TF
regulatory evidence in the KG.

\begin{center}\rule{0.5\linewidth}{0.5pt}\end{center}

\section{What TF Enrichment Data is
Available}\label{what-tf-enrichment-data-is-available}

TF enrichment is performed using \textbf{gseapy (enrich)} against
curated TF--target databases. Results are stored as edges from TF gene
nodes to cohort or network nodes:

{\def\LTcaptype{none} % do not increment counter
\begin{longtable}[]{@{}
  >{\raggedright\arraybackslash}p{(\linewidth - 2\tabcolsep) * \real{0.5000}}
  >{\raggedright\arraybackslash}p{(\linewidth - 2\tabcolsep) * \real{0.5000}}@{}}
\toprule\noalign{}
\begin{minipage}[b]{\linewidth}\raggedright
Edge type
\end{minipage} & \begin{minipage}[b]{\linewidth}\raggedright
Meaning
\end{minipage} \\
\midrule\noalign{}
\endhead
\bottomrule\noalign{}
\endlastfoot
\texttt{tf\_regulates\_cohort\_degs} & TF whose targets are enriched in
a specific cohort's DEG list \\
\texttt{tf\_regulates\_network\_genes} & TF whose targets are enriched
in a WGCNA module \\
\end{longtable}
}

Key properties on TF edges: \texttt{padj} (adjusted p-value),
\texttt{overlap\_genes} (list of shared target genes),
\texttt{n\_targets} (total targets regulated).

Available for: \textbf{Study 1 cohorts} and \textbf{WGCNA modules} (gene
layer only).

\begin{center}\rule{0.5\linewidth}{0.5pt}\end{center}

\section{Querying TF Enrichment}\label{querying-tf-enrichment}

\subsection{TFs Regulating High-Producer
DEGs}\label{tfs-regulating-high-producer-degs}

Which transcription factors regulate genes upregulated in high vs.~low
Nivo producers in Study 1?

Show TFs with adj. p-value \textless{} 0.01 whose targets are enriched
among upregulated DEGs in C144 high producers.

\subsection{TFs Linked to a Specific
Module}\label{tfs-linked-to-a-specific-module}

Which TFs regulate genes in WGCNA module M1 (pro-productivity)?

What transcription factors control the anti-productivity gene module
M18?

\subsection{Tracing a Gene to Its Upstream
Regulators}\label{tracing-a-gene-to-its-upstream-regulators}

What transcription factors regulate HSPA5 in the context of
high-producer cohorts?

Trace gene KCTD14 upstream to its transcription factors.

\begin{center}\rule{0.5\linewidth}{0.5pt}\end{center}

\section{Master Regulator
Identification}\label{master-regulator-identification}

A \textbf{master regulator} is a TF that controls a large number of
target genes across multiple modules or cohorts.

Which transcription factors have the broadest target gene scope across
all pro-productivity cohorts in Study 1?

Key master regulators identified across pro-productivity (M1/M7/M8) and
anti-productivity (M16/M18) modules:

{\def\LTcaptype{none} % do not increment counter
\begin{longtable}[]{@{}lll@{}}
\toprule\noalign{}
TF & Module association & padj (representative) \\
\midrule\noalign{}
\endhead
\bottomrule\noalign{}
\endlastfoot
\textbf{E2F4} & Anti-productivity (M18) & 3.7e−122 \\
FOXM1 & Anti-productivity & Significant \\
E2F1 & Anti-productivity & Significant \\
MYC & Anti-productivity & Significant \\
NFE2L2 (NRF2) & Pro-productivity (M1) & Significant \\
FOXO1 & Pro-productivity (M1) & Significant \\
PPARA & Pro-productivity (M1) & Significant \\
FOXA1 & Pro-productivity (M1) & Significant \\
SMARCA4 & Pro-productivity & Significant \\
KLF5 & Pro-productivity & Significant \\
\end{longtable}
}

\begin{tcolorbox}[enhanced jigsaw, arc=.35mm, bottomrule=.15mm, bottomtitle=1mm, breakable, colback=white, colbacktitle=quarto-callout-note-color!10!white, colframe=quarto-callout-note-color-frame, coltitle=black, left=2mm, leftrule=.75mm, opacityback=0, opacitybacktitle=0.6, rightrule=.15mm, title=\textcolor{quarto-callout-note-color}{\faInfo}\hspace{0.5em}{Note}, titlerule=0mm, toprule=.15mm, toptitle=1mm]

\textbf{E2F4} is the most statistically significant TF in the dataset
(padj = 3.7×10⁻¹²²), regulating genes in the anti-productivity cell
cycle/DNA replication module M18. Targeting E2F4 or its upstream pathway
could suppress cell proliferation programs and redirect resources toward
productivity.

\end{tcolorbox}

\begin{center}\rule{0.5\linewidth}{0.5pt}\end{center}

\section{Checking if a TF is Itself Differentially
Expressed}\label{checking-if-a-tf-is-itself-differentially-expressed}

Is NFE2L2 (NRF2) differentially expressed in high vs.~low producers in
Study 1?

Is FOXM1 a member of any WGCNA module? What is its gs\_qp score?

Does E2F4 have a significant MOFA Factor 1 weight?

\begin{center}\rule{0.5\linewidth}{0.5pt}\end{center}

\section{Combining TF Evidence with Other
Layers}\label{combining-tf-evidence-with-other-layers}

For the highest-confidence regulatory targets:

Show TFs that (1) regulate upregulated DEGs in high producers, (2) are
themselves upregulated in high producers, AND (3) are members of a
pro-productivity WGCNA module.

\bookmarksetup{startatroot}

\chapter{Hypothesis Testing}\label{hypothesis-testing}

Interrogating the KG when you arrive with a prior

\hfill\break

\section{Overview}\label{sec-hypothesis}

This chapter is for users who arrive at the KG with a \textbf{prior}---a
gene of interest from a paper, a pathway hypothesis, or a specific
biological mechanism they want to test. The KG lets you rapidly check
how much evidence exists for any prior across all analytical layers.

\begin{center}\rule{0.5\linewidth}{0.5pt}\end{center}

\section{I Have a Gene --- What Does the KG
Know?}\label{i-have-a-gene-what-does-the-kg-know}

Use this multi-layer gene audit as your default starting point whenever
you have a gene of interest.

Tell me everything the KG knows about gene HSPA5: differential
expression, WGCNA module, MOFA weight, phenotype correlations, and
published evidence.

The KG will return evidence from:

\begin{enumerate}
\def\labelenumi{\arabic{enumi}.}
\tightlist
\item
  {DEG/DAA} Is it DE in high vs.~low producers? In how many cohorts? In
  which direction?
\item
  {WGCNA} Which module does it belong to? What is its kME? Is the module
  correlated with qP/titre/IVCD?
\item
  {MOFA} Does it have a significant Factor 1 weight? Positive or
  negative?
\item
  Spearman correlation with qP, titre, and IVCD directly.
\item
  {Literature} Has it been manipulated in CHO cells? With what effect?
\end{enumerate}

\begin{center}\rule{0.5\linewidth}{0.5pt}\end{center}

\section{I Have a Pathway --- Which Genes in it Are
Productivity-Associated?}\label{i-have-a-pathway-which-genes-in-it-are-productivity-associated}

Which genes in the ``Protein processing in endoplasmic reticulum'' KEGG
pathway are differentially expressed in BMS Study 1 high vs.~low
producer cohorts?

Show all genes in the glycolysis / gluconeogenesis pathway with positive
Spearman correlation to qP in Study 1.

\begin{center}\rule{0.5\linewidth}{0.5pt}\end{center}

\section{I Have a Gene List --- Check Against All Evidence
Layers}\label{i-have-a-gene-list-check-against-all-evidence-layers}

I have a list of genes: {[}HSPA5, DNAJB9, MERTK, FICD, ZCCHC24{]}. For
each, show: (1) DEG status in high vs.~low Nivo producers, (2) WGCNA
module, (3) MOFA Factor 1 weight, (4) literature evidence.

\begin{center}\rule{0.5\linewidth}{0.5pt}\end{center}

\section{I Have a Phenotype --- What Genes Are Most
Correlated?}\label{i-have-a-phenotype-what-genes-are-most-correlated}

Which genes have the strongest positive Spearman correlation with
specific productivity (qP) in Study 1?

Rank proteins by their Spearman correlation with titre in Study 1,
descending.

Which metabolites are most negatively correlated with qP?

\begin{center}\rule{0.5\linewidth}{0.5pt}\end{center}

\section{Testing Specific Biological
Hypotheses}\label{testing-specific-biological-hypotheses}

\subsection{ER Stress / Unfolded Protein Response (UPR)
Hypothesis}\label{er-stress-unfolded-protein-response-upr-hypothesis}

\emph{Prior:} High producers activate UPR to expand ER folding capacity.

Show all genes involved in the UPR or ER stress response that are
upregulated in high producers in Study 1.

Is the ``protein processing in endoplasmic reticulum'' KEGG pathway
enriched in the pro-productivity WGCNA module M1?

Which UPR-related transcription factors (ATF4, ATF6, XBP1) are
differentially expressed in high producers?

\textbf{KG verdict:} ✅ Strongly supported. Module M1 is enriched for ER
processing / UPR. HSPA5, DNAJB9, and FICD carry positive MOFA Factor 1
weights. Literature evidence (PMID 32624757, 24081924) corroborates
GRP78/BiP overexpression as a productivity lever.

\begin{center}\rule{0.5\linewidth}{0.5pt}\end{center}

\subsection{Secretory Pathway
Hypothesis}\label{secretory-pathway-hypothesis}

\emph{Prior:} High producers have enhanced protein secretion machinery.

Are genes in the ``protein export'' and ``antigen processing'' pathways
upregulated in high producers in Study 1?

Show the top co-expressed genes with HSPA5 in WGCNA module M1.

\textbf{KG verdict:} ✅ Supported. Protein export and antigen processing
are both enriched in WGCNA M1 (pro-productivity module).

\begin{center}\rule{0.5\linewidth}{0.5pt}\end{center}

\subsection{Metabolic Bottleneck Hypothesis (Lactate / Ammonia /
ROS)}\label{metabolic-bottleneck-hypothesis-lactate-ammonia-ros}

\emph{Prior:} Runaway lactate accumulation limits productivity in low
producers.

Which metabolites significantly differ between high and low producers in
Study 1? Highlight lactate.

Show genes in the lactate metabolism pathway (KEGG) with negative
Spearman correlation to titre.

Are there published interventions (gene KO/OE) that reduce lactate
accumulation in CHO cells?

\begin{center}\rule{0.5\linewidth}{0.5pt}\end{center}

\subsection{Chaperone / Difficult-to-Express (DTE)
Hypothesis}\label{chaperone-difficult-to-express-dte-hypothesis}

\emph{Prior:} DTE proteins benefit from co-expression of specific
chaperones.

Show all chaperone genes (HSP family, DNAJ family) upregulated in high
producers in Study 1.

What publications link chaperone overexpression to improved productivity
of difficult-to-express proteins?

\begin{center}\rule{0.5\linewidth}{0.5pt}\end{center}

\subsection{Apoptosis-Resistance
Hypothesis}\label{apoptosis-resistance-hypothesis}

\emph{Prior:} Blocking apoptosis extends culture longevity and increases
total titre.

Are anti-apoptotic genes (BCL2, BCL-XL family) upregulated in high
producers in any BMS cohort?

Which published studies link apoptosis-resistance interventions
(combined knockouts) to reduced lactate and improved productivity?

\bookmarksetup{startatroot}

\chapter{Target Prioritisation}\label{target-prioritisation}

Shortlisting high-confidence genetic engineering candidates

\hfill\break

\section{Overview}\label{sec-targets}

This chapter is the \textbf{deliverable section}---where multi-layer
evidence converges into a ranked shortlist of overexpression (OE) and
knockout (KO) candidates for experimental validation.

\begin{center}\rule{0.5\linewidth}{0.5pt}\end{center}

\section{What Makes a High-Confidence
Candidate?}\label{what-makes-a-high-confidence-candidate}

The KG scores candidates across six dimensions:

\begin{center}\rule{0.5\linewidth}{0.5pt}\end{center}

\section{Overexpression Candidates}\label{overexpression-candidates}

Overexpression candidates are genes that, when expressed at higher
levels, are expected to \textbf{increase} productivity.

\textbf{Criteria:}

\begin{itemize}
\tightlist
\item
  Upregulated in high producers in ≥ 7/10 cohorts
\item
  Positive Spearman correlation with qP (r \textgreater{} 0.5)
\item
  Positive MOFA Factor 1 weight
\item
  Member of pro-productivity WGCNA module (e.g., M1)
\item
  Published OE evidence showing positive productivity effect (where
  available)
\end{itemize}

Identify overexpression candidates: genes upregulated in high producers
across ≥7/10 cohorts, with positive MOFA Factor 1 weight and positive qP
correlation \textgreater{} 0.5.

Rank overexpression candidates by number of supporting high-producer
cohorts, then by abs(log2FC).

\subsection{Example: HSPA5}\label{example-hspa5}

{DEG} {WGCNA} {MOFA} {Literature}

{\def\LTcaptype{none} % do not increment counter
\begin{longtable}[]{@{}
  >{\raggedright\arraybackslash}p{(\linewidth - 2\tabcolsep) * \real{0.5000}}
  >{\raggedright\arraybackslash}p{(\linewidth - 2\tabcolsep) * \real{0.5000}}@{}}
\toprule\noalign{}
\begin{minipage}[b]{\linewidth}\raggedright
Evidence layer
\end{minipage} & \begin{minipage}[b]{\linewidth}\raggedright
Value
\end{minipage} \\
\midrule\noalign{}
\endhead
\bottomrule\noalign{}
\endlastfoot
DEG consistency & Up in 7/10 cohorts \\
Spearman qP & r = +0.52 \\
Spearman titre & r = +0.38 \\
MOFA Factor 1 weight & +0.392 \\
WGCNA module & M1 (kME = 0.660); ER processing / UPR enriched \\
TF regulators & NFE2L2 (NRF2), FOXO1, PPARA, FOXA1 \\
Literature & PMIDs 32624757, 24081924, 1373378 --- GRP78/BiP OE → ↑
titre, ↑ qP \\
\textbf{Recommendation} & ▲ \textbf{OVEREXPRESS} \\
\end{longtable}
}

\begin{center}\rule{0.5\linewidth}{0.5pt}\end{center}

\section{Knockout Candidates}\label{knockout-candidates}

Knockout candidates are genes that, when eliminated, are expected to
\textbf{relieve} a productivity bottleneck.

\textbf{Criteria:}

\begin{itemize}
\tightlist
\item
  Downregulated in high producers in ≥ 7/10 cohorts (or upregulated,
  reflecting a gene that suppresses productivity)
\item
  Negative Spearman correlation with qP (r \textless{} −0.5)
\item
  Negative MOFA Factor 1 weight
\item
  Member of anti-productivity WGCNA module (e.g., M18)
\item
  Published KO evidence showing positive productivity effect (where
  available)
\end{itemize}

Identify knockout candidates: genes with negative MOFA Factor 1 weight,
downregulated in ≥7/10 high-producer cohorts, and negative Spearman qP
correlation.

\subsection{Example: ST14}\label{example-st14}

{DEG} {WGCNA} {MOFA} {Literature}

{\def\LTcaptype{none} % do not increment counter
\begin{longtable}[]{@{}
  >{\raggedright\arraybackslash}p{(\linewidth - 2\tabcolsep) * \real{0.5000}}
  >{\raggedright\arraybackslash}p{(\linewidth - 2\tabcolsep) * \real{0.5000}}@{}}
\toprule\noalign{}
\begin{minipage}[b]{\linewidth}\raggedright
Evidence layer
\end{minipage} & \begin{minipage}[b]{\linewidth}\raggedright
Value
\end{minipage} \\
\midrule\noalign{}
\endhead
\bottomrule\noalign{}
\endlastfoot
DEG consistency & Down in 10/10 cohorts \\
Spearman qP & r = −0.51 \\
Spearman titre & r = −0.51 \\
MOFA Factor 1 weight & −0.303 (anti-productivity loading) \\
WGCNA module & M18 (kME = 0.595); Cell cycle / DNA replication
enriched \\
TF regulators & E2F4 (padj = 3.7×10⁻¹²²), FOXM1, E2F1, MYC \\
Literature & PMID 29777593 --- ST14 KO → ↓ fragmentation, ↑ product
purity \\
\textbf{Recommendation} & ▼ \textbf{KNOCKOUT} \\
\end{longtable}
}

\begin{center}\rule{0.5\linewidth}{0.5pt}\end{center}

\section{Multi-Omics Concordance
Filter}\label{multi-omics-concordance-filter}

Apply concordance filtering to further refine candidates:

From the 63-gene shortlist, which have concordant RNA--protein
regulation (mRNA up and protein up for OE candidates, or mRNA down and
protein down for KO candidates)?

\begin{center}\rule{0.5\linewidth}{0.5pt}\end{center}

\section{Molecule-Specific vs.~Shared
Targets}\label{molecule-specific-vs.-shared-targets}

Compare overexpression candidates between Nivo and C144. Which targets
are shared across both molecules, and which are molecule-specific?

Show knockout candidates that appear in both Nivo and C144 high-vs-low
cohorts.

Shared targets are preferred for \textbf{cell-line-wide} engineering;
molecule-specific targets may be relevant for product-specific quality
or expression challenges.

\begin{center}\rule{0.5\linewidth}{0.5pt}\end{center}

\section{Exporting the Shortlist}\label{exporting-the-shortlist}

Export the top 20 overexpression candidates as a ranked table with
columns: gene symbol, DEG cohort count, Spearman qP, MOFA Factor 1
weight, WGCNA module, literature PMID.

The exported table provides a ready-to-use input for experimental
planning and prioritisation discussions.

\bookmarksetup{startatroot}

\chapter{KG Utilities \& Navigation}\label{kg-utilities-navigation}

Searching, cross-referencing identifiers, filtering, and exporting

\hfill\break

\section{Overview}\label{sec-utilities}

This chapter covers the practical mechanics of navigating the KG:
searching by identifier, cross-referencing databases, filtering results,
and exporting data.

\begin{center}\rule{0.5\linewidth}{0.5pt}\end{center}

\section{Searching by Identifier}\label{searching-by-identifier}

The KG supports search by multiple identifier types:

{\def\LTcaptype{none} % do not increment counter
\begin{longtable}[]{@{}ll@{}}
\toprule\noalign{}
Entity & Supported identifiers \\
\midrule\noalign{}
\endhead
\bottomrule\noalign{}
\endlastfoot
Gene & Gene symbol (e.g., \texttt{HSPA5}), NCBI Gene ID \\
Protein & UniProt accession (e.g., \texttt{P11021}) \\
Metabolite & ChEBI ID, metabolite name \\
Pathway & KEGG pathway ID (e.g., \texttt{hsa04141}), GO term ID \\
\end{longtable}
}

Find gene HSPA5. Show its NCBI Gene ID, UniProt accession of its encoded
protein, and all associated KEGG pathways.

Search for metabolite lactate using its ChEBI ID CHEBI:16651.

\begin{center}\rule{0.5\linewidth}{0.5pt}\end{center}

\section{Cross-Referencing
Identifiers}\label{cross-referencing-identifiers}

What is the UniProt accession for the protein encoded by HSPA5?

What CHO gene symbol corresponds to NCBI Gene ID 14828?

What is the KEGG ortholog for human LDHA in CHO cells?

\begin{center}\rule{0.5\linewidth}{0.5pt}\end{center}

\section{Finding Orthologs}\label{finding-orthologs}

The KG maps CHO genes to human and mouse orthologs via KEGG Orthologs
(KO terms).

Find the human ortholog of CHO gene LOC100760815.

Which human genes are orthologs of the top 10 hub genes in WGCNA module
M1?

This is particularly useful when interpreting CHO findings in light of
human biology or when searching for mechanistic context in human
databases.

\begin{center}\rule{0.5\linewidth}{0.5pt}\end{center}

\section{Filtering Results}\label{filtering-results}

\subsection{By Study}\label{by-study}

Show DEG results for HSPA5 restricted to BMS internal studies only
(exclude public omics data).

Show only public omics datasets with DEG results for KCTD14.

\subsection{By Molecule}\label{by-molecule}

Filter WGCNA module enrichment results to Nivo only.

Show MOFA Factor 1 weights for genes in the C144 dataset only.

\subsection{By Statistical Threshold}\label{by-statistical-threshold}

Recommended defaults:

{\def\LTcaptype{none} % do not increment counter
\begin{longtable}[]{@{}ll@{}}
\toprule\noalign{}
Threshold & Value \\
\midrule\noalign{}
\endhead
\bottomrule\noalign{}
\endlastfoot
DEG adj. p-value & \textless{} 0.05 \\
DEG \textbar log2FC\textbar{} & \textgreater{} 2 \\
Spearman correlation & \textbar r\textbar{} \textgreater{} 0.5 \\
MOFA weight (top percentile) & \textgreater{} 80th percentile \\
WGCNA module membership & kME \textgreater{} 0.5 \\
\end{longtable}
}

Show DEGs with adj. p-value \textless{} 0.01 and
\textbar log2FC\textbar{} \textgreater{} 3 in the Nivo high vs.~low
cohort.

\begin{center}\rule{0.5\linewidth}{0.5pt}\end{center}

\section{Understanding Sample
Coverage}\label{understanding-sample-coverage}

Not all omics layers are available for all samples:

{\def\LTcaptype{none} % do not increment counter
\begin{longtable}[]{@{}lll@{}}
\toprule\noalign{}
Omics layer & Study 1 days & Study 2 days \\
\midrule\noalign{}
\endhead
\bottomrule\noalign{}
\endlastfoot
Transcriptomics & 3, 5, 9, 11, 14 & --- \\
Proteomics & 3, 5, 9, 11, 14 & 5, 7, 9, 11 \\
Metabolomics & 3, 5, 9, 11, 14 & --- \\
\end{longtable}
}

Which samples in Study 1 have all three omics layers measured on the
same day?

\begin{center}\rule{0.5\linewidth}{0.5pt}\end{center}

\section{Exporting Query Results}\label{exporting-query-results}

Results can be exported as:

\begin{itemize}
\tightlist
\item
  \textbf{CSV / TSV} --- gene lists, ranked tables, enrichment results
\item
  \textbf{JSON} --- for programmatic downstream processing
\item
  \textbf{Network formats} --- for import into Cytoscape or other graph
  tools
\end{itemize}

Export all WGCNA module M1 members with kME, gs\_qp, and gs\_titre as a
CSV.

Export the top 50 genes by positive MOFA Factor 1 weight across all
three omics layers.

\begin{center}\rule{0.5\linewidth}{0.5pt}\end{center}

\section{Tips for Effective NLQ
Queries}\label{tips-for-effective-nlq-queries}

\begin{enumerate}
\def\labelenumi{\arabic{enumi}.}
\tightlist
\item
  \textbf{Be specific about the entity type} --- ``gene'', ``protein'',
  ``metabolite''
\item
  \textbf{Name the cohort explicitly} --- include molecule and study
\item
  \textbf{State directionality} --- ``upregulated'', ``positively
  correlated'', ``with negative MOFA weight''
\item
  \textbf{Include thresholds} --- exact numbers rather than vague
  qualifiers (``strong'', ``significant'')
\item
  \textbf{Layer your filters progressively} --- start broad, then add
  constraints
\item
  \textbf{Reference specific modules or factors by name} --- ``WGCNA
  module M1'', ``MOFA Factor 1''
\end{enumerate}

\begin{tcolorbox}[enhanced jigsaw, arc=.35mm, bottomrule=.15mm, bottomtitle=1mm, breakable, colback=white, colbacktitle=quarto-callout-tip-color!10!white, colframe=quarto-callout-tip-color-frame, coltitle=black, left=2mm, leftrule=.75mm, opacityback=0, opacitybacktitle=0.6, rightrule=.15mm, title=\textcolor{quarto-callout-tip-color}{\faLightbulb}\hspace{0.5em}{Tip}, titlerule=0mm, toprule=.15mm, toptitle=1mm]

If a query returns too many results, add a threshold filter. If it
returns nothing, relax one constraint at a time (e.g., lower the cohort
consistency requirement from 7/10 to 5/10).

\end{tcolorbox}

\bookmarksetup{startatroot}

\chapter{Glossary}\label{glossary}

\section{Key Terms}\label{key-terms}

\begin{description}
\tightlist
\item[\textbf{adj. p-value (adjusted p-value)}]
P-value corrected for multiple testing using the Benjamini-Hochberg (BH)
method. Values below 0.05 are considered statistically significant in
this KG.
\item[\textbf{C144}]
One of the two BMS antibody molecules for which in-house multi-omics
data is available in Study 1.
\item[\textbf{CHO (Chinese Hamster Ovary)}]
The cell line used for biopharmaceutical protein production. All BMS
in-house data in this KG comes from CHO cells.
\item[\textbf{Cohort}]
A pairwise comparison group used for differential expression analysis
(e.g., Nivo high vs.~low producers, or clone A day 5 vs.~clone B day 9).
The KG contains 883 cohorts.
\item[\textbf{DAA (Differential Abundance Analysis)}]
Statistical analysis identifying proteins or metabolites that differ
significantly in abundance between two conditions. Method: T-test with
time as covariate.
\item[\textbf{DEG (Differentially Expressed Gene)}]
A gene whose mRNA expression differs significantly between two
conditions. Method: DESeq2 with time as covariate.
\item[\textbf{DTE (Difficult-to-Express)}]
Proteins with complex folding or disulfide-bond requirements that make
them hard to produce at high levels in CHO cells.
\item[\textbf{Factor (MOFA Factor)}]
A latent variable identified by MOFA that captures coordinated variation
across multiple omics layers. Factor 1 is the most strongly correlated
with productivity phenotypes (qP +0.848, IVCD +0.860, titre +0.904).
\item[\textbf{GEM (Genome-scale Metabolic Model)}]
A constraint-based computational model of all metabolic reactions in a
cell. Used for in-silico gene perturbation to predict effects on
secretion flux. \emph{GEM integration is in progress.}
\item[\textbf{GO (Gene Ontology)}]
A standardised vocabulary of biological terms. Sub-ontologies: BP
(Biological Process), MF (Molecular Function), CC (Cellular Component).
\item[\textbf{GS (Gene Significance)}]
In WGCNA, the Pearson correlation between a gene's expression and a
trait. Stored as \texttt{gs\_qp}, \texttt{gs\_ivcd}, and
\texttt{gs\_titre}.
\item[\textbf{GSEA (Gene Set Enrichment Analysis)}]
Enrichment method that uses a ranked gene list to detect coordinated
shifts in pathway activity. Reports Normalised Enrichment Score (NES).
\item[\textbf{IVCD (Integral Viable Cell Density)}]
The area under the viable cell density--time curve. Measures cumulative
biomass produced during a culture run.
\item[\textbf{KEGG (Kyoto Encyclopedia of Genes and Genomes)}]
A database of biological pathways, reactions, enzymes, and compounds.
Used for pathway enrichment and metabolic network traversal in this KG.
\item[\textbf{kME (Module Eigengene Connectivity)}]
In WGCNA, the correlation between a gene's expression and its module
eigengene. Higher kME (closer to 1) indicates the gene is a better
representative (hub gene) of the module.
\item[\textbf{KO (Knockout)}]
Complete elimination of a gene's function, typically via CRISPR-Cas9.
Contrast with KD (knockdown, partial reduction) and OE (overexpression).
\item[\textbf{ME (Module Eigengene)}]
The first principal component of all features within a WGCNA module.
Represents the overall expression pattern of the module across samples.
\item[\textbf{MOFA (Multi-Omics Factor Analysis)}]
An unsupervised dimensionality-reduction framework that integrates
multiple omics layers into shared latent factors.
\item[\textbf{Nivo (Nivolumab)}]
One of the two BMS antibody molecules (anti-PD-1 mAb) for which in-house
data is available in Studies 1 and 2.
\item[\textbf{NES (Normalised Enrichment Score)}]
The GSEA enrichment score normalised for gene set size and the number of
permutations. Positive NES = enrichment in the first group; negative NES
= enrichment in the second group.
\item[\textbf{NLQ (Natural Language Query)}]
A plain-English question typed into the KG interface, which is
automatically converted to a Cypher graph query.
\item[\textbf{OE (Overexpression)}]
Engineering a cell to produce more of a specific protein, typically via
stable transgene insertion.
\item[\textbf{ORA (Over-Representation Analysis)}]
Enrichment method that tests whether a binary gene list contains more
pathway members than expected by chance.
\item[\textbf{qP (Specific Productivity)}]
The rate of product secretion per cell per unit time (pg/cell/day). A
key productivity phenotype in the KG.
\item[\textbf{Titre}]
Total product concentration in the culture supernatant at harvest
(mg/L). Reflects both qP and viable cell density over time.
\item[\textbf{TF (Transcription Factor)}]
A protein that binds DNA regulatory elements to activate or repress gene
expression. TF enrichment identifies master regulators of DEG programs.
\item[\textbf{TOM (Topological Overlap Matrix)}]
A measure of network interconnectedness used in WGCNA. Genes with high
TOM are densely connected and tend to cluster into the same module.
\item[\textbf{TPM (Transcripts Per Million)}]
A normalisation unit for RNA-seq expression data. Accounts for
sequencing depth and transcript length.
\item[\textbf{UPR (Unfolded Protein Response)}]
A cellular stress response activated when misfolded proteins accumulate
in the ER. Associated with pro-productivity gene expression programs in
high CHO producers.
\item[\textbf{WGCNA (Weighted Gene Co-expression Network Analysis)}]
A network-based method that groups co-expressed genes, proteins, or
metabolites into modules and correlates module activity with biological
traits.
\end{description}


\backmatter


\end{document}
